\documentclass{article}
\usepackage{graphicx}
\usepackage{booktabs}
\usepackage{float}
\usepackage{amsmath}
\usepackage{hyperref}

\title{Analysis of Bias in Matryoshka vs. Standard Embeddings}
\author{Experimental Report}
\date{\today}

\begin{document}

\maketitle

\section{Introduction}
This report analyzes the bias characteristics of Matryoshka Representation Learning (MRL) embeddings compared to standard non-MRL embeddings. We compare two models:
\begin{itemize}
    \item \textbf{MRL Model}: \texttt{tomaarsen/mpnet-base-nli-matryoshka}
    \item \textbf{Baseline (Non-MRL)}: \texttt{tomaarsen/mpnet-base-nli}
\end{itemize}

We investigate four research questions (RQ1--RQ4) focusing on bias subspace survival, concentration, and debiasing effectiveness at various prefix dimensions ($d \in \{64, 128, \dots, 768\}$).

\section{RQ1: Bias Subspace Survival}
\textit{Does the bias subspace determined at full dimension $D$ align with the subspace at prefix dimension $d$?}

We measure the subspace alignment error between the projected full bias subspace $\Pi_d(B^{(D)})$ and the local prefix bias subspace $B^{(d)}$. Lower error indicates better survival.

\begin{table}[H]
    \centering
    \caption{Subspace Alignment Error (lower is better)}
    \begin{tabular}{c c c}
    \toprule
    Dimension ($d$) & Baseline Error & MRL Error \\
    \midrule
    64 & 0.289 & \textbf{0.284} \\
    128 & 0.262 & \textbf{0.259} \\
    256 & 0.211 & \textbf{0.208} \\
    384 & 0.160 & \textbf{0.159} \\
    512 & 0.100 & \textbf{0.098} \\
    \bottomrule
    \end{tabular}
\end{table}

\textbf{Analysis:} The MRL model consistently shows slightly lower alignment error across all dimensions. This indicates that the bias direction in MRL models is more stable and "survives" truncation better than in standard models, likely due to the nested structure of MRL training.

\section{RQ3: Bias Concentration}
\textit{Does MRL concentrate bias information in early dimensions?}

We calculate the concentration score $C(d) = \text{Var}(b^{(d)}) / \text{Var}(b^{(D)})$, representing the proportion of total bias captured in the first $d$ dimensions.

\begin{table}[H]
    \centering
    \caption{Bias Concentration $C(d)$ (higher indicates earlier concentration)}
    \begin{tabular}{c c c}
    \toprule
    Dimension ($d$) & Baseline $C(d)$ & MRL $C(d)$ \\
    \midrule
    64 & 0.085 & \textbf{0.158} \\
    128 & 0.165 & \textbf{0.241} \\
    256 & 0.287 & \textbf{0.376} \\
    384 & 0.497 & \textbf{0.558} \\
    512 & 0.695 & \textbf{0.719} \\
    \bottomrule
    \end{tabular}
\end{table}

\textbf{Analysis:} The MRL model shows significantly higher bias concentration in early dimensions. At $d=64$, MRL captures 15.8\% of the total bias variance compared to only 8.5\% for the baseline. At $d=128$, this gap widens (24.1\% vs 16.5\%). This confirms that MRL packs concept information (including bias) into the prefix.

\section{RQ2: Debiasing at Prefix Level}
\textit{Does debiasing at the full dimension $D$ effectively remove bias at prefix dimension $d$?}

We measure the residual bias (L2 norm of projection onto bias subspace) at each prefix dimension after applying DeepSoftDebias at $D=768$.

\begin{table}[H]
    \centering
    \caption{Residual Bias after Full-Dimension Debiasing}
    \begin{tabular}{c c c}
    \toprule
    Dimension ($d$) & Baseline Residual & MRL Residual \\
    \midrule
    64 & \textbf{0.106} & 0.150 \\
    128 & 0.139 & \textbf{0.115} \\
    256 & \textbf{0.110} & 0.144 \\
    384 & 0.171 & \textbf{0.169} \\
    \midrule
    768 (Full) & 0.252 & \textbf{0.223} \\
    \bottomrule
    \end{tabular}
\end{table}

\textbf{Analysis:}
\begin{itemize}
    \item \textbf{At Full Dimension ($D=768$):} MRL is more amenable to global debiasing (Residual 0.223 vs 0.252).
    \item \textbf{At Prefix ($d=64$):} MRL retains \textit{more} residual bias (0.150) than the baseline (0.106), likely due to the high concentration of bias features in the first 64 dimensions that global debiasing fails to remove.
    \item \textbf{At Prefix ($d=128$):} MRL actually retains \textit{less} residual bias (0.115 vs 0.139), showing that by dimension 128, the global debiasing operation effectively cleans the subspace.
\end{itemize}

\section{RQ4: Adaptive Lambda}
\textit{Does the optimal debiasing parameter $\lambda$ change with dimension?}

We performed a grid search for $\lambda_2$ (trade-off parameter) to minimize bias while maintaining similarity $>0.9$.

\begin{itemize}
    \item \textbf{MRL Model at $d=64$:} Optimal $\lambda \approx 0.8$.
    \item \textbf{Baseline Model at $d=64$:} Optimal $\lambda \approx 0.5$.
\end{itemize}

\textbf{Conclusion:} Because MRL concentrates bias in the early dimensions (as seen in RQ3), a stronger debiasing penalty ($\lambda=0.8$) is required to effectively remove it compared to the baseline model.

\section{Discussion: Intrinsic vs. Downstream Bias Correlation}
\textit{How does intrinsic embedding bias correlate with downstream task bias under hierarchical compression?}

We compared the Intrinsic Bias Concentration $C(d)$ (from RQ3) with the Downstream Stereotype Score (SS) on the StereoSet dataset.

\begin{itemize}
    \item \textbf{General Trend:} Positive correlation. Higher bias concentration $C(d)$ tends to result in higher downstream stereotyping scores.
    \item \textbf{Model Comparison at $d=128$:}
    \begin{itemize}
        \item \textbf{MRL:} High Concentration ($C \approx 0.24$) $\rightarrow$ \textbf{Higher SS (54.5)}
        \item \textbf{Baseline:} Lower Concentration ($C \approx 0.16$) $\rightarrow$ \textbf{Lower SS (53.7)}
    \end{itemize}
\end{itemize}

The results indicate that the "packed" bias in MRL \textit{does} translate to slightly higher stereotyping scores in downstream tasks compared to the baseline. The high density of bias information in the prefix ($d=64, 128$) makes the MRL model more prone to exhibiting stereotypes when truncated, compared to the baseline which has less information (and less bias) in those early dimensions.

\end{document}
